% !TeX root = ../../../geos_mpm_manual.tex
\section{Robustness controls: deletion, deformation-gradient reset, and particle splitting}
\label{sec:robustness-controls}
\index{robustness controls}
\index{particle deletion}
\index{deformation gradient!reset}
\index{particle splitting}

This section documents solver-side mechanisms that deliberately alter the active particle set or the particle kinematic state to keep large-deformation MPM calculations well posed.  These controls are not substitutes for mesh/time-step convergence studies or physically calibrated constitutive regularization, but they are useful guardrails for impact, fragmentation, contact, pore collapse, and highly damaged calculations.  The controls described here act after the ordinary explicit transfer/update sequence described in Section~\ref{sec:solver-steps} and are related to, but distinct from, the contact and cohesive-zone algorithms described in Sections~\ref{sec:contact-options} and~\ref{sec:cohesive-zone-implementation}.  User-facing parameters are listed in the generated solver attribute appendix and, when they are written by ParticleFileWriter, in the PFW controls of Chapter~\ref{chap:pfw}.

\subsection{Particle deletion and active-particle compaction}
\label{subsec:robust-particle-deletion}
\index{particleDeleteFlag}
\index{maxParticleJacobian}
\index{minParticleJacobian}
\index{maxParticleVelocity}
\index{constitutiveUpdateFlag}

GEOS-MPM uses a deferred deletion model.  During the step, several kernels or events set a per-particle integer flag,
\begin{equation}
  \texttt{particleDeleteFlag}_p = 1,
\end{equation}
when particle $p$ should no longer participate in subsequent active-particle loops.  The physical erase is performed later by \texttt{deleteBadParticles}, which moves particle-field wrappers to host memory, builds a set of indices with the delete flag set, erases those entries from the particle subregion, and reconstructs the active-particle index set.  Deleting is therefore a topology-changing operation; it removes the particle mass, history variables, constitutive state, and any diagnostic fields attached to the particle.

The principal automatic deletion criteria are:
\begin{enumerate}
  \item \textbf{Constitutive failure flag.}  If a material wrapper exposes \texttt{constitutiveUpdateFlag}, \texttt{updateSolverDependencies} copies it to the solver field \texttt{particleConstitutiveUpdateFlag}.  A negative value at the first material point flags the particle for deletion.  This path lets a constitutive model request deletion after an unrecoverable update, for example failed return mapping or a material-specific erosion condition.
  \item \textbf{Deformation-gradient determinant limits.}  During \texttt{particleKinematicUpdate}, the solver computes
  \begin{equation}
    J_p = \det \mathbf{F}_p.
  \end{equation}
  The particle is flagged if
  \begin{equation}
    J_p \le J_{\min}\quad\hbox{or}\quad J_p \ge J_{\max},
    \label{eq:robust-j-deletion}
  \end{equation}
  where $J_{\min}=\texttt{minParticleJacobian}$ and $J_{\max}=\texttt{maxParticleJacobian}$.  In the reviewed source snapshot the constructor defaults are $J_{\min}=0.1$ and $J_{\max}=10$.  This criterion removes particles that have inverted, collapsed, or expanded beyond the range for which the mapped particle volume and constitutive response are trustworthy.
  \item \textbf{Velocity overflow and maximum velocity.}  The kinematic update also checks whether any component is so large that squaring it would overflow a floating-point value.  Such particles are flagged before kinetic-energy or speed calculations can fail.  The registered user parameter \texttt{maxParticleVelocity} sets the intended speed limit through \texttt{maxParticleVelocitySquared}.  In the reviewed code, the overflow guard is active; the speed-squared accumulation for the ordinary maximum-speed comparison appears inside the overflow branch after a \texttt{break}, so that ordinary \texttt{maxParticleVelocity} deletion should be verified before relying on it as an active erosion rule.
  \item \textbf{Out-of-domain particles.}  At the end-of-step cleanup stage, \texttt{flagOutOfRangeParticles} removes particles that cannot be mapped safely on the next step.  For \texttt{SinglePoint} particles the center must remain inside the global domain plus one buffer-cell layer, excluding periodic directions.  For \texttt{SinglePointBSpline} particles the center must remain inside the physical global domain because the cubic B-spline basis already needs the available buffer support.  For \texttt{CPDI} particles, all eight current-domain corners generated from the center and three $\mathbf{r}$ vectors must remain in range.  Periodic directions are skipped because particle centers are wrapped during repartitioning.
  \item \textbf{Event-driven machining.}  The \texttt{MachineSample} event sets delete flags for particles outside an idealized dogbone, cylinder, or Brazilian-disk sample geometry.  This is a constructive-geometry deletion mechanism rather than a numerical-failure guard.
\end{enumerate}

A compact view of the deletion path is:
\begin{lstlisting}[language=Python,caption={Deferred particle deletion in GEOS-MPM.}]
for particle p:
    if constitutiveUpdateFlag[p,0] < 0:
        particleDeleteFlag[p] = 1
    J = det(F[p])
    if J <= minParticleJacobian or J >= maxParticleJacobian:
        particleDeleteFlag[p] = 1
    if any velocity component would overflow when squared:
        particleDeleteFlag[p] = 1

# Later, after grid resize/domain cleanup:
for particle p:
    if particle center or CPDI corner cannot map to available nodes:
        particleDeleteFlag[p] = 1

erase all particles with particleDeleteFlag == 1
rebuild activeParticleIndices
\end{lstlisting}

Because particle deletion is non-conservative unless the removed material has negligible mass or represents intentional erosion, production studies should report the deletion thresholds, the number or mass of deleted particles, and whether deletion occurs at outflow boundaries, from constitutive failure, or from numerical pathologies.

\subsection{Deformation-gradient reset}
\label{subsec:robust-defgrad-reset}
\index{domainResetType}
\index{VP-Hencky deformation-gradient reset}
\index{resetDefGradForFullyDamagedParticles}
\index{resetDefGradForMeltedParticles}
\index{resetDefGradForScaledSurfaceParticles}
\index{ResetDeformationGradient}
\index{fluid material!deformation-gradient reset}

The deformation-gradient reset is a physical-state reparameterization for particles whose constitutive response no longer requires the accumulated finite deviatoric stretch stored in $\mathbf{F}$.  It is intended for fluid-like materials, melted material, fully damaged material, and explicitly requested scaled surface particles.  The independently integrated particle-volume ratio
\begin{equation}
  J_p=\frac{V_p}{V_p^0}
  \label{eq:reset-volume-ratio}
\end{equation}
is the target Jacobian; $\det\mathbf{F}_p$ is retained as a geometric validity check.  The input \texttt{domainResetType} selects \texttt{isotropicPolar}, which is the default and reproduces the previous behavior, or \texttt{vpHencky}.

\paragraph{Isotropic-polar reset.}
For \texttt{isotropicPolar}, the solver extracts the polar rotation $\mathbf{Q}_p$ from the old deformation gradient and replaces the stretch by
\begin{equation}
  \mathbf{U}^{\mathrm{iso}}_p =
  \begin{cases}
    \operatorname{diag}(J_p^{1/2},J_p^{1/2},1), & \texttt{planeStrain}=1,\\[3pt]
    J_p^{1/3}\mathbf{I}, & \texttt{planeStrain}=0,
  \end{cases}
\end{equation}
so that
\begin{equation}
  \mathbf{F}^{\mathrm{new}}_p=\mathbf{Q}_p\mathbf{U}^{\mathrm{iso}}_p.
  \label{eq:isotropic-polar-reset}
\end{equation}
This removes all active-dimensional deviatoric stretch while preserving rotation and $J_p$.

\paragraph{VP-Hencky reset.}
For \texttt{vpHencky}, the solver first forms the raw current domain matrix from $\mathbf{F}^{\mathrm{old}}_p$ and the reference particle vectors as in Equation~\eqref{eq:cpdi-raw-domain-from-f}.  The requested half-domain measure is
\begin{equation}
  M_p^{\star}=J_p M_p^0,
  \label{eq:reset-target-measure}
\end{equation}
where $M_p^0$ is the reference half-domain area in plane strain or volume in three dimensions.  The projection of Equations~\eqref{eq:vp-hencky-domain-diagonal}--\eqref{eq:vp-hencky-path} retains the largest feasible fraction of the old deviatoric Hencky stretch subject to
\begin{equation}
  D_p < \texttt{neighborRadius}.
\end{equation}
The new physical deformation gradient is recovered from the projected domain using
\begin{equation}
  \mathbf{R}^{\mathrm{new}}_p
  =\mathbf{R}^0_p(\mathbf{F}^{\mathrm{new}}_p)^T,
  \qquad
  \mathbf{F}^{\mathrm{new}}_p
  =\left[(\mathbf{R}^0_p)^{-1}\mathbf{R}^{\mathrm{new}}_p\right]^T,
  \label{eq:reset-f-from-domains}
\end{equation}
with particle half-edge vectors stored row-wise in $\mathbf{R}$.  If the old domain already satisfies the diagonal constraint and its measure agrees with Equation~\eqref{eq:reset-target-measure}, the VP-Hencky reset leaves $\mathbf{F}$ unchanged.  If the fixed measure and diagonal bound are incompatible, the reset preserves $M_p^{\star}$ and uses the equal-measure $\beta=0$ domain; particle splitting is then required to make the bound feasible.  Singular or invalid domain geometry falls back to \texttt{isotropicPolar}.

The activation routes are:
\begin{description}
  \item[Fluid materials.]  Subregions whose constitutive catalog name is \texttt{Gas} or \texttt{Liquid} are eligible every cleanup pass.
  \item[Melted particles.]  If \texttt{resetDefGradForMeltedParticles=1}, particles with \texttt{particleMeltFlag > 0.5} are eligible.
  \item[Fully damaged particles.]  If \texttt{resetDefGradForFullyDamagedParticles=1}, particles with \texttt{particleDamage > 0.9999999} are eligible.  CPDI domain scaling is not required for this selection.
  \item[Scaled surface particles and events.]  If \texttt{resetDefGradForScaledSurfaceParticles=1}, scaled CPDI particles with surface-like flags are eligible.  The \texttt{ResetDeformationGradient} MPM event sets this one-pass request, and the solver clears it after the cleanup pass.
\end{description}

Changing $\mathbf{F}$ without changing its associated reference data would rotate or distort surfaces, material axes, cohesive kinematics, and particle domains on the next update.  The reset therefore rebases the stored reference quantities so their current values are preserved.  In particular,
\begin{align}
  \boldsymbol{n}^{s,0}_p
    &\leftarrow \operatorname{cof}(\mathbf{F}^{\mathrm{new}}_p)^{-1}
      \boldsymbol{n}^{s}_p,
  &
  \boldsymbol{x}^{s,0}_p
    &\leftarrow (\mathbf{F}^{\mathrm{new}}_p)^{-1}
      \boldsymbol{x}^{s}_p,\\
  \boldsymbol{t}^{s,0}_p
    &\leftarrow \operatorname{cof}(\mathbf{F}^{\mathrm{new}}_p)^{-1}
      \boldsymbol{t}^{s}_p.
\end{align}
Normals and directional material vectors are normalized after rebasing; traction magnitude is retained.  Fiber-like material directions use $(\mathbf{F}^{\mathrm{new}})^{-1}$, while area-like material directions use the inverse cofactor, matching their ordinary forward updates.  For cohesive particles the reference cohesive deformation gradient is changed so that $\mathbf{F}\mathbf{F}_{\mathrm{cz}}^{-1}$ is unchanged.  Current particle $\mathbf{r}$ vectors are regenerated from the new $\mathbf{F}$ for every reset route, and the stored $\dot{\mathbf{F}}$ and velocity gradient are cleared because the reset is not a physical rate increment.

\begin{lstlisting}[language=Python,caption={Deformation-gradient reset selection and dispatch.}]
for active particle p:
    selected = (Gas or Liquid material
                or melted particle selected by input
                or fully damaged particle selected by input
                or one-pass scaled-surface selection)
    if not selected:
        continue
    J = particleVolume[p] / particleReferenceVolume[p]
    if domainResetType == "isotropicPolar":
        Fnew = polar_rotation(Fold) @ isotropic_stretch(J)
    else:  # vpHencky
        raw_domain = Fold @ reference_domain
        target_measure = J * reference_domain_measure
        projected_domain = vp_hencky_project(raw_domain,
                                               target_measure,
                                               neighborRadius)
        Fnew = recover_F(reference_domain, projected_domain)
    rebase surface, material-direction, and cohesive reference data
    regenerate current particle-domain vectors from Fnew
resetDefGradForScaledSurfaceParticles = 0
\end{lstlisting}

\subsection{Particle splitting}
\label{subsec:robust-particle-splitting}
\index{subdivideParticles}
\index{oversizedParticleTreatment}
\index{oversizedParticleResetRankParticleCountThreshold}
\index{defGradResetMaxParticleDomainToGridCellRatio}
\index{resetDefGradForOversizedParticles}
\index{particleSubdivideFlag}
\index{particleCopyFlag}
\index{CPDI!particle splitting}
\index{VP-Hencky feasibility}

Particle splitting is a conservative topology change used when the requested particle-domain measure and the neighbor-radius diagonal constraint cannot both be satisfied.  It is called during topology preparation, after events and before particle maps, ghosts, and active-particle indices are refreshed.  The input is tri-state:
\begin{description}
  \item[\texttt{subdivideParticles=-1}] use the conditional default;
  \item[\texttt{subdivideParticles=0}] disable subdivision;
  \item[\texttt{subdivideParticles=1}] enable subdivision explicitly.
\end{description}
The constructor default is $-1$.  Automatic subdivision resolves to enabled only when \texttt{cpdiDomainScaling=1}, \texttt{cpdiDomainScalingType="vpHencky"}, and at least one linear \texttt{CPDI} particle subregion is present.  It remains disabled by default for single-point B-spline particles and for deformation-gradient reset alone.  Explicitly setting the value to 1 enables the VP-Hencky reset feasibility test for B-spline or other particle types that carry domain vectors.  In a mixed-particle simulation, the automatic policy acts only on the linear CPDI subregions that caused it to be enabled.

Oversized-domain handling is selected separately with \texttt{oversizedParticleTreatment}.  The legacy \texttt{resetDefGradForOversizedParticles=1} path is equivalent to \texttt{oversizedParticleTreatment="resetDeformationGradient"} when the new option is left at \texttt{"none"}.  With \texttt{oversizedParticleTreatment="split"}, a particle whose full domain diagonal exceeds \texttt{defGradResetMaxParticleDomainToGridCellRatio} times the active grid-cell diagonal is routed through the same conservative child-creation path used below.  The split pass first applies the VP-Hencky fixed-volume feasibility count when relevant, then adds longest domain directions until one bisection pass would bring the oversized diagonal below the configured limit, or all active directions have been bisected.  Deformation-gradient reset is used as a count-control fallback only on ranks whose total local particle count exceeds \texttt{oversizedParticleResetRankParticleCountThreshold}; a negative threshold disables that fallback.

For either VP-Hencky path, define
\begin{equation}
  M_{\max}=\left(\frac{0.99999\,\texttt{neighborRadius}}
                            {2\sqrt{d}}\right)^d.
\end{equation}
The CPDI test uses the raw unscaled domain measure.  The reset test uses the independently integrated target $J_pM_p^0$ for particles selected for reset.  If the required measure $M_p$ exceeds $M_{\max}$, the solver selects the smallest number $k\le d$ such that
\begin{equation}
  \frac{M_p}{2^k}\le M_{\max},
  \label{eq:split-feasibility-count}
\end{equation}
and splits along the $k$ longest current domain vectors.  A pass can bisect each active direction at most once; a domain with $M_p/M_{\max}>2^d$ is split again on a subsequent explicit step.  CPDI integration-domain scaling may temporarily relax its integration measure during this transition, while a VP-Hencky physical reset always preserves particle volume.

If $k$ directions are selected, the parent particle is replaced by $2^k$ children: one child reuses the parent slot and $2^k-1$ particles are appended.  Mass, current volume, and reference volume are divided by $2^k$.  The selected current and reference half-edge vectors are halved, and current and reference centers are shifted by all sign combinations of the halved vectors:
\begin{align}
  m_{p_i}&=\frac{m_p}{2^k},
  &V_{p_i}&=\frac{V_p}{2^k},
  &V^0_{p_i}&=\frac{V^0_p}{2^k},\\
  \boldsymbol{x}_{p_i}
    &=\boldsymbol{x}_p+
      \sum_{a\in\mathcal{S}_p}s_{i,a}\boldsymbol{r}_{p_i,a},
  &s_{i,a}&\in\{-1,+1\}.
\end{align}
All remaining particle fields and constitutive history fields are copied from the parent through \texttt{particleCopyFlag}.  Since the reference and current vectors are bisected consistently, every child retains the parent's deformation gradient before any subsequent reset.  Surface-position fields are stored relative to the particle center; their current and reference values are translated by the corresponding child-center offsets so that every child initially represents the same absolute surface point as the parent.

Particles with an active cohesive-zone state are not subdivided by this generic path.  Their cached reference-node maps, shape functions, cohesive reference geometry, and constitutive history cannot be conservatively partitioned by merely copying the parent state.  For such particles, CPDI scaling uses the exact feasible integration-measure cap when necessary, while an infeasible physical VP-Hencky reset preserves particle volume and reports that the equal-volume domain remains over the diagonal limit.

For explicit subdivision in configurations that select neither VP-Hencky path, the solver retains the historical length-based criterion, which bisects every active domain vector longer than $0.49999$ times the smallest active grid spacing.  This preserves backward compatibility for inputs that already set \texttt{subdivideParticles=1}.

\begin{lstlisting}[language=Python,caption={VP-Hencky feasibility-driven particle splitting.}]
Mmax = (0.99999 * neighborRadius / (2 * sqrt(active_dimensions)))**active_dimensions
for particle p selected by an enabled VP-Hencky path:
    Mrequired = raw_CPDI_measure_or_reset_target_measure(p)
    if Mrequired <= Mmax:
        continue
    choose the fewest longest domain directions needed so Mrequired / 2**k <= Mmax
    create 2**k children
    divide mass, current volume, and reference volume by 2**k
    halve current and reference domain vectors in the selected directions
    shift current and reference centers to tile the parent domain
    copy remaining particle and constitutive fields
rebuild activeParticleIndices
\end{lstlisting}

Splitting conserves parent mass and volume but changes the particle quadrature and can increase particle count rapidly.  Production studies should report whether the conditional or explicit policy was used and monitor the emitted subdivision count.  A persistent need for repeated splitting may indicate that the neighbor radius, grid resolution, particle density, or physical volume expansion requires reconsideration.
